% paper.tex — your ACM acmart skeleton, started in Week 1.
% Copy this into a new Overleaf project (Appendix A walks you through it),
% or compile locally with: pdflatex paper.tex
%
% acmart.cls is v2.18 (2026/06/01) as of mid-2026 — re-check the literal
% version string in Overleaf's package list before you rely on a feature.
% For IEEE instead, swap to: \documentclass[conference]{IEEEtran}
%
% Submit single-column for review; switch to [sigconf] for camera-ready.
% Add the "anonymous" option for double-blind venues -> "ANONYMOUS AUTHOR(S)".
\documentclass[manuscript,review,anonymous]{acmart}

\begin{document}

\title{A Working Title: How <AI method> Advances <your CS domain>}

\author{Anonymous Author(s)}   % double-blind: do NOT put your name yet

\begin{abstract}
% Kent Beck's four sentences — write these LAST, but stub them now:
% (1) the problem  (2) why it is a problem  (3) the contribution
% (4) the implication if it works.
TODO: one paragraph. Problem. Why it matters. What we do. Why it matters that we did it.
\end{abstract}

\maketitle

\section{Introduction}
% Swales CARS: establish territory -> establish the niche (the gap) -> occupy it.
TODO: What is the area? What gap did you find (cite the absence)? What do you claim?

\section{Related Work}
% Whom do you build on, and whom do you beat? This grows into your Lit Review (Week 4).
TODO

\section{Method}
% What you did, reproducibly. Another team should be able to rerun it from this text.
TODO

\section{Results}
% What you found. Baselines, multiple seeds, effect sizes + CIs --- not one lucky run.
TODO

\section{Discussion}
% What it means. Limitations. Threats to validity. Do not over-claim past the data.
TODO

\section{Future Work}
TODO

\begin{acks}
% AI-disclosure goes HERE for ACM if you used generative AI substantively.
% Sample (ACM): "ChatGPT was utilized to generate sections of text..."
% Assistive grammar/spell-check needs no disclosure; generated text/code/ideas do.
TODO: disclose any substantive AI assistance per the venue policy (Appendix C).
\end{acks}

\bibliographystyle{ACM-Reference-Format}
\bibliography{references}   % your Zotero-exported references.bib (Week 3)

\end{document}
